\subsubsection{Objetivos específicos [Pendientes: 1. acordar métricas y umbrales, 2. precisar tamaño de biblioteca, 3. confirmar plazos de cumplimiento]}\label{sec:objetivos}
Los cuatro objetivos del póster se conservan como base y se ajustan en su redacción al alcance vigente de software. La retroalimentación del PAFP solicita reforzar ese enfoque, pero no contiene una lista definitiva que sustituya los objetivos previos~\cite{pafp2026,poster2026}. Por esta razón, se propone la siguiente formulación de trabajo, cuya medición se concretará con los criterios de las fases y los valores pendientes indicados en este título.
\begin{enumerate}
\item \textbf{OE1.} Diseñar e implementar un mecanismo de búsqueda que recupere escenarios de la biblioteca a partir de la ubicación o el perímetro observado y de condiciones meteorológicas actuales obtenidas mediante una API, conservando la identificación de los casos seleccionados para comprobar la correspondencia entre la consulta y sus resultados.
\item \textbf{OE2.} Construir una biblioteca consultable de simulaciones WFDS precalculadas para las igniciones y condiciones acordadas del sector de estudio, mediante ejecuciones en HPC y las modalidades de paralelización que se verifiquen como compatibles, con metadatos que permitan identificar entradas, configuración y salidas de cada caso.
\item \textbf{OE3.} Configurar WFDS con niveles de complejidad de parámetros pertinentes para representar la propagación en las coberturas urbanas y forestales presentes en el sector, preparando entradas reproducibles y documentando sus restricciones, de manera que sus resultados puedan incorporarse al flujo de la aplicación.
\item \textbf{OE4.} Desarrollar e integrar la aplicación de escritorio con animaciones georreferenciadas del avance estimado del fuego y actualización de consultas, para que sus resultados y limitaciones puedan ser interpretados por el usuario previsto sin formación especializada en modelado, mediante tareas de evaluación que aporten evidencia de esa comprensión.
\end{enumerate}
